%% kapa: exam=1 round=29 year=2018 period=2 form=A subject=appraisal_law subject_ko=감정평가관계법규 seq=19 number=19
%% source: Q-Net 감정평가사 기출문제 (cst00309) artlSeq=5202808 file=extracted/1008/5202808/2018년 제29회 감정평가사 1차시험 기출문제/2018년 제29회 감정평가사 1차시험 기출문제/2교시 A형.pdf
%% conversion: kordoc 4.24.51 -> tools/gen1.py (auto); review=proofread-sonnet; flags=none
\begin{question}{19}
부동산 가격공시에 관한 법령상 지가의 공시 등에 관한 설명으로 옳지 않은 것은?
\begin{choices}
  \item 표준지로 선정되어 개별공시지가를 결정ㆍ공시하지 아니하는 토지의 경우 해당 토지의 표준지공시지가를 개별공시지가로 본다.
  \item 개별공시지가 조사ㆍ산정의 기준에는 지가형성에 영향을 미치는 토지 특성조사에 관한 사항이 포함되어야 한다.
  \item 관계공무원등이 표준지가격의 조사ㆍ평가를 위해 택지에 출입하고자 할 때에는 점유자를 알 수 없거나 부득이한 사유가 있는 경우를 제외하고는 출입할 날의 3일 전에 그 점유자에게 일시와 장소를 통지하여야 한다.
  \item 표준지공시지가의 공시기준일은 1월 1일이며, 일부 지역을 지정하여 해당 지역에 대한 공시기준일을 따로 정할 수는 없다.
  \item 개별공시지가의 결정ㆍ공시에 소요되는 비용은 대통령령으로 정하는 바에 따라 그 일부를 국고에서 보조할 수 있다.
\end{choices}
\end{question}
\answer{4}
